\section{Related Inverse Problems and Algorithmic Context}\label{sec:related}
\subsection{Representation, sharing and regularization}
Positive Laplace inversion is a first-kind Fredholm problem in which small singular values limit stable recovery \cite{discrete,hansen}. The main methodological distinction is the unknown being estimated. CONTIN and UPEN recover distributions stabilized by smoothing, whereas PALMA combines entropy and sparsity \cite{contin,upen,palma}. Their penalties express different preferences for admissible solutions; none makes finite observations identify an unrestricted positive distribution. The nonnegative smoothing framework of Butler, Reeds and Dawson supplies an earlier reference \cite{brd}.

SCORE uses shared discrete exponential components, while LOCODOSY fits local spectral windows whose orders can differ \cite{score,locodosy}. These are particularly relevant comparators for the atomic RAI-S and DOME-S models. Global sharing pools evidence across resonances; local fitting limits how much unrelated spectral structure is coupled. Neither choice removes the need to determine which components contribute at each chemical-shift position. OUTSCORE provides a further component-resolution approach, and GNAT makes these distinctions explicit in a common processing environment \cite{outscore,gnat}. Multivariate curve resolution and nonnegative matrix factorization provide broader precedents for shared-factor models \cite{mcr,nmf}. In TRAIn-MF the attenuation factor is additionally constrained to $KS$, with a known kernel $K$. Normalizing the profiles removes a scale freedom but not all factorization ambiguity; a profile can remain broad or multimodal.

Table~\ref{tab:externalmethods} contrasts seven relevant external methods across distributional, component-based and learned representations. It identifies feasible comparisons and their contracts. The frozen primary benchmark is supplemented by the executed SCORE, LOCODOSY and PALMA-objective comparisons in Appendix~\ref{app:externalcomparison}; the remaining entries are context rather than executed rankings.
\begin{table}[htbp]
\caption{Properties and requirements for external comparisons. Availability of source code or trained weights is distinct from validation on a common dataset.}\label{tab:externalmethods}
\centering\small
\begin{tabularx}{\textwidth}{>{\raggedright\arraybackslash}p{0.16\textwidth}*{2}{>{\raggedright\arraybackslash}X}}
\toprule
Method & Representation and stabilizing information & Requirement for a fair DOSY comparison\\
\midrule
CONTIN \cite{contin} & Nonnegative distribution; constrained smoothing & Match kernel, grid convention and regularization selection; compare prediction and distributional error.\\
PALMA \cite{palma} & Positive distribution; entropy and sparsity & Fix the penalty mixture, noise constraint and stopping rule without test-data tuning.\\
UPEN \cite{upen} & Distribution; locally adapted penalty & Use a verified one-dimensional Laplace implementation, not a second inversion over chemical shift.\\
SCORE \cite{score} & Shared discrete components and spectral amplitudes & Declare order selection, initialization, amplitude constraints and computation budget.\\
LOCODOSY \cite{locodosy} & Local multivariate fits and local component orders & Derive windows and orders from training data; distinguish local contributions from global species.\\
DRECT \cite{drect} & Learned Laplace-NMR reconstruction & Match acquisition sampling, units, normalization and training domain; test distribution shift.\\
DREAM \cite{dream} & Learned reconstruction with uncertainty estimates & Match the same input contracts and assess uncertainty calibration on independent data.\\
\bottomrule
\end{tabularx}
\end{table}

Code is available for these methods or related implementations, and pretrained DOSY checkpoints are available for DRECT and DREAM. This makes direct comparisons feasible, but does not make their interfaces interchangeable. Distributional methods do not natively estimate an atomic order; LOCODOSY's local component count is not a global species count; and a learned model may require sampling or normalization different from the present protocol. Adapters must preserve physical units and the observations used for evaluation. In particular, two-dimensional relaxation tools \cite{separable2d,upen2d,mupen} cannot treat the chemical-shift axis as a second Laplace dimension. The primary frozen benchmark does not establish a ranking against these seven external methods. The additional executed comparisons concern the explicitly attributed implementations, adapters, parameter rules and endpoints in Appendix~\ref{app:externalcomparison}; CONTIN, UPEN, DRECT and DREAM remain unexecuted in that common protocol.

\subsection{Optimization, identification and learning}
Efficient optimization does not overcome a lack of information. FLINT addresses the computational cost of distribution estimation, and FISTA supplies convergence results for specified convex objectives \cite{flint,fista}; those results do not automatically cover a bilinear factorization or adaptive support and order changes. ESPIRA estimates exponential parameters through rational approximation and structured matrices \cite{espira}. Its sampling assumptions, like the assumptions of any identification theorem, must be checked for the model under consideration \cite{teicher}. Noiseless uniqueness in a bounded-order atomic class and stable recovery of nearby rates in noise are different claims.

DRECT and DREAM also obtain information from their training distributions \cite{drect,dream}. Their evaluation therefore requires independent test cases and checks under distribution shift. RAI-Net has a different role: it changes the order of refinement proposals in an adaptive solver. Neither the term ``learning'' nor use of a physical kernel makes it equivalent to a physics-informed neural network, and the present selector's conditional Gaussian calculations do not provide DREAM's learned uncertainty model or a calibrated chemical confidence interval.

The distinction between representation, optimization and evidence extends beyond Laplace NMR. Fractional-diffusion studies compare optimizers for recovering derivative orders and boundary conditions \cite{fractionalmeta}; hybrid searches have been investigated for reconstruction from incomplete tomographic measurements \cite{hybridct}. Studies of physics-based neural networks further motivate sensitivity analyses across domain geometry and training hyperparameters \cite{pinngeometry}. These examples involve different forward operators and unknowns: a fractional partial differential equation, tomographic projection data, and differential-equation solutions, respectively. They provide methodological context, not direct DOSY performance comparators or validation of RAI-Net.
